% Installation on the Samsung YP-S3. Summary of "For the owner: install, go
% back, rescue" in YPS3-PORT.md, which has the exact commands.
\subsection{Installing Rockbox on the \playertype}

\warn{The \playertype{} port is experimental. Rockbox Utility does not support
it and there are no builds for it on the Rockbox website: you build Rockbox
from the source tree that contains this port. Not all of the steps below
have been tried on a player yet; the steps marked \emph{[NOT YET TESTED]} in
\fname{YPS3-PORT.md} are unverified.}

The exact procedure, with all commands, is in the file \fname{YPS3-PORT.md}
in the top folder of the source tree, in the section ``For the owner:
install, go back, rescue''. This section summarises it. All commands run
from the top of the source tree on a Linux PC with Python~3, pyusb,
dosfstools and gio.

\subsubsection{How it works}
\begin{itemize}
  \item The \daps{} own bootloader looks for a file called
    \fname{S3.ROM} in the top folder of the music drive every time it is
    switched on. If it finds one, it checks it, writes it as the firmware,
    deletes it and starts it. Installing Rockbox and going back to the
    original firmware both work this way. The \playerman{} bootloader itself
    is never changed.
  \item The scripts \fname{yps3-install.sh} and \fname{yps3-restore.sh} in
    \fname{utils/s5l8700\_dfu/} only check an image and copy it to the drive
    as \fname{S3.ROM}. Nothing on the \dap{} changes until it
    is switched off and on.
  \item The image that installs Rockbox contains both the original firmware
    and the Rockbox bootloader. After the install, the \dap{} starts Rockbox.
    For the original firmware, keep the \ButtonPower{} slider pushed through
    the whole start animation; when it stops on its last frame, let go (see
    \reference{ref:Dualboot}).
  \item DFU mode is built into the \daps{} chip and always works, whatever is
    on the flash. It lets the PC start a program on the \dap{} over USB without
    writing anything (see \reference{ref:yps3_rescue}).
\end{itemize}

\subsubsection{What you need}
\begin{description}
  \item[A backup of the \daps{} flash.] It is made with
    \fname{utils/s5l8700\_dfu/yps3\_backup.py} and kept in
    \fname{nanddump/backup/}. The file \fname{hidden\_area\_ecc.img} in it is
    the only copy of the original firmware, so keep a copy of the directory
    somewhere other than your PC.
  \item[The builds.] The Rockbox bootloader (\fname{rockboot.ebn}, in the
    build directory \fname{build-yps3}) and the main Rockbox build
    (\fname{rockbox.zip}, made with \fname{make zip} in
    \fname{build-yps3-full}).
  \item[The images.] \fname{utils/s5l8700\_dfu/yps3\_mkrom.py extract}
    takes the original firmware out of the backup as
    \fname{S3-original.ROM}; it checks the image against the checksum stored
    in it. \fname{yps3\_mkrom.py combine} puts it together with
    \fname{rockboot.ebn} as \fname{S3-combined.ROM}, the image that installs
    Rockbox. An official \playerman{} \fname{S3.ROM} for the \playertype{}
    can replace \fname{S3-original.ROM}. \fname{yps3\_mkrom.py info FILE}
    shows what an image is.
\end{description}

\subsubsection{The music drive}
When Rockbox is connected to a computer by USB, it shows its USB screen and
the \dap{} appears as a drive: the partition of about 3.7~GB on a 4~GB
\dap{}. This is the music drive that the scripts are given. As a safety net
against picking another disk, they refuse anything under 512~MB. The
original firmware shows no drive, only an MTP device. Before Rockbox is
installed, get the music drive by starting Rockbox over USB in DFU mode (see
\reference{ref:yps3_rescue}).

\subsubsection{Installation}
Charge the battery first. While the \dap{} writes a new firmware, keep the
cable unplugged and don't switch it off until a firmware has started.
\begin{enumerate}
  \item Once, before installing, write the original firmware over itself:
    run \fname{yps3-restore.sh MUSIC\_DRIVE}, then eject
    the drive, unplug the cable, switch the \dap{} off (hold \ButtonPower{}
    until the screen goes dark) and on. The \dap{} shows a short upgrade
    screen, the original firmware starts, and \fname{S3.ROM} is gone from the
    drive. Nothing changes, but this proves that the \dap{} takes
    \fname{S3.ROM} from the drive, which is also the way back.
  \item Get the music drive again and run
    \fname{yps3-install.sh MUSIC\_DRIVE}. It checks
    \fname{S3-combined.ROM}, unpacks \fname{rockbox.zip} to the drive
    (the \fname{.rockbox} directory) and copies the image to the drive as
    \fname{S3.ROM}.
  \item Eject the drive, unplug the cable, switch the \dap{} off and on. The
    \dap{} writes the new firmware, \fname{S3.ROM} disappears and Rockbox
    starts. With Rockbox installed the start animation takes longer than
    before, because the \playerman{} bootloader loads the whole 25~MB image
    instead of 9~MB.
\end{enumerate}

\subsubsection{When \fname{S3.ROM} stays}
If there was no upgrade screen and \fname{S3.ROM} is still on the drive after
switching off and on:
\begin{itemize}
  \item make sure the \dap{} was really off and the cable unplugged;
  \item check that the file is in the top folder of the music drive, named
    exactly \fname{S3.ROM}, about 25~MB, with no \fname{S3J.ROM} next to it;
  \item \emph{[NOT YET TESTED]} the \playerman{} bootloader may only look for
    updates after the original firmware has been connected to a computer
    once: start it, plug it in, unplug, and try again;
  \item otherwise stop: nothing has changed and the \dap{} works as before.
    \fname{S3.ROM} can stay or be deleted.
\end{itemize}

\subsubsection{\label{ref:yps3_rescue}Rescue: nothing starts (DFU mode)}
To enter DFU mode:
\begin{enumerate}
  \item Plug the USB cable in.
  \item Hold the \ButtonPower{} slider down and keep holding it.
  \item Press a pin into the reset hole for 1--2 seconds, then release the pin.
  \item Wait a moment, then release \ButtonPower{}. If both are released
    almost together, the \dap{} starts its firmware normally instead; try
    again and wait a little longer before letting go of \ButtonPower{}.
\end{enumerate}
The screen stays dark. \fname{s5l8700\_dfu.py list} (in \fname{utils/s5l8700\_dfu/})
must show \fname{0419:0145 Samsung YP-S3}. Then:
\begin{description}
  \item[Rockbox on the drive is fine.]
    \fname{s5l8700\_dfu.py run build-yps3/rockboot.ebn} starts the Rockbox
    bootloader, which starts Rockbox from the drive.
  \item[Rockbox on the drive is broken too.] Start the bootloader from
    \fname{build-yps3-dev} the same way, then load \fname{rockbox.yps3}
    from the PC with \fname{yps3\_nanddump.py}. See
    \fname{YPS3-PORT.md} for the commands.
\end{description}
Then go back to the original firmware (see \reference{ref:yps3_uninstall}) or
install again. Starting a program over DFU writes nothing; switching off
brings the \dap{} back to how it was.

If the music drive can't be mounted or shows errors, \fname{YPS3-PORT.md}
describes how to repair it with \fname{fsck.vfat} and, only as a last resort,
how to reformat it, which erases all music and \fname{.rockbox} on the
\dap{}.
